%% =========================================================
%% Chapter: Matrix Product Operator Symmetries and Non-Invertible
%% Symmetries.
%% The basic notions of symmetries given by matrix product operators:
%% on-site and matrix product unitary group symmetries, group operators
%% built from a three-cocycle and their anomaly, fusion algebras and
%% their Perron--Frobenius dimensions, non-invertible symmetries
%% (Fibonacci, Ising, Kramers--Wannier), symmetric matrix product states
%% with their nonnegative integer representations, and the zipper
%% condition. The virtual-level theory (fusion and action tensors,
%% gauge, associator, obstruction) is Chapter ch:mpo_algebras.
%% Lean modules: TNLean/MPS/Symmetry/MPOSymmetry/{Defs,NIMRep,Character,
%%   Examples}.lean (Sections sec:mposym_fusion and sec:mposym_symmetric).
%% Planned Lean module for the Perron--Frobenius dimension:
%%   TNLean/MPS/Symmetry/MPOSymmetry/Dimension.lean.
%% **Scope restriction (periodic boundary):** Garre-Rubio, Lootens and
%%   Molnar (arXiv:2203.12563) define the algebra by closedness under
%%   arbitrary boundary conditions (Papers/2203.12563/REsubmission.tex,
%%   lines 328--331) and symmetry by invariance of the arbitrary-boundary
%%   subspace (lines 431--434); Sections sec:mposym_fusion and
%%   sec:mposym_symmetric use only the periodic-boundary consequences
%%   (lines 361--362 and 567--568). Documented in
%%   docs/paper-gaps/glm23_mpo_symmetric_mps_scope.tex.
%% =========================================================

\chapter{Matrix Product Operator Symmetries and Non-Invertible Symmetries}%
\label{ch:mpo_symmetry_basics}

A symmetry of a spin chain is usually a group acting site by site. Two
generalizations appear in one-dimensional lattice models. The symmetry
operators may be matrix product operators rather than tensor products of
one-site unitaries, and they may multiply according to fusion rules
$O_aO_b=\sum_cN_{ab}^cO_c$ rather than a group law, so that some of them
have no inverse
\cite{Bultinck2017Anyons,GarreRubio2023MPOSym,Cirac2021Matrix}. This
chapter introduces both generalizations for a reader who knows matrix
product states. Section~\ref{sec:mposym_invertible} treats invertible
symmetries: on-site actions, groups represented by matrix product unitaries,
and group operators built from a three-cocycle, whose cohomology class
obstructs an on-site form. Section~\ref{sec:mposym_fusion} defines fusion
algebras of matrix product operators and their Perron--Frobenius dimensions,
and Section~\ref{sec:mposym_noninvertible} presents the non-invertible
examples: Fibonacci, Ising, and Kramers--Wannier duality.
Section~\ref{sec:mposym_symmetric} derives the integrality consequences for
symmetric matrix product states, and Section~\ref{sec:mposym_zipper} states
the local zipper condition. The virtual-level theory, with fusion and action
tensors, their gauge freedom, the associator, and the cohomological
obstruction, is Chapter~\ref{ch:mpo_algebras}.

Throughout, $T_a$ denotes a matrix product operator tensor
(Definition~\ref{def:mpo_tensor}) of bond dimension $\chi_a$, and
$O_a^{(N)}$ its periodic operator on $N$ sites
(Definition~\ref{def:mpo_family}),
\begin{align}
    O_a^{(N)}
     & =\sum_{i_1,\ldots,i_N}\sum_{j_1,\ldots,j_N}
    \tr(T_a^{i_1j_1}\cdots T_a^{i_Nj_N})
    \ket{i_1,\ldots,i_N}\bra{j_1,\ldots,j_N},
    \label{eq:mposym_periodic_operator}
\end{align}
where the first physical index is the output. We write $\ket{V^{(N)}(A)}$ for
the periodic vector of a matrix product state tensor $A$. For two operator
tensors, $T_aT_b$ denotes the stacked tensor
$(T_aT_b)^{ij}=\sum_mT_a^{im}\otimes T_b^{mj}$ of bond dimension
$\chi_a\chi_b$, whose periodic operator is $O_a^{(N)}O_b^{(N)}$, and
$T_a\cdot A$ denotes the matrix product state tensor
$(T_a\cdot A)^i=\sum_jT_a^{ij}\otimes A^j$, whose periodic vector is
$O_a^{(N)}\ket{V^{(N)}(A)}$. Chain lengths are positive, $N\geq1$.

%% ---------------------------------------------------------
\section{Invertible symmetries}%
\label{sec:mposym_invertible}

\begin{remark}[On-site symmetries]%
    \label{rem:mposym_onsite}
    \uses{def:on_site_symmetric, thm:virtual_symmetry_eq,
        cor:virtual_rep_cocycle, def:mpo_family}
    A representation $u\colon G\to\GL_d(\C)$ acts on a chain of $N$ sites by
    $u_g^{\otimes N}$. This is the periodic operator of the bond-dimension-one
    tensor $T_g^{ij}=u(g)_{ij}$, so on-site symmetries are the matrix product
    operator symmetries of bond dimension one
    \cite[Section~4]{GarreRubio2023MPOSym}. Invariance of a matrix product
    state under them is Definition~\ref{def:on_site_symmetric}. For a normal
    tensor the physical action is carried to the bond by
    Theorem~\ref{thm:virtual_symmetry_eq}, and the bond operators form a
    projective representation whose class in $H^2(G,\C^\times)$ is the
    invariant of Corollary~\ref{cor:virtual_rep_cocycle}.
    % Source anchor: arXiv:2203.12563, Papers/2203.12563/REsubmission.tex,
    %   lines 658 (on-site operators as bond-dimension-one blocks) and
    %   840--856 (reduction to the second-cohomology classification).
\end{remark}

\begin{remark}[Groups represented by matrix product unitaries]%
    \label{rem:mposym_mpu_groups}
    \uses{def:mpu, def:mpug_positive_mpu, def:mpug_group_representation}
    A symmetry need not act site by site to be local. A matrix product
    unitary (Definition~\ref{def:mpu}) is a matrix product operator whose
    periodic operator is unitary at every length $N>1$. A finite group is
    represented by simple injective matrix product unitaries when every
    $U_g$ is injective, simple and unitary on every nonempty chain
    (Definition~\ref{def:mpug_positive_mpu}), and
    $U_g^{(N)}U_h^{(N)}=U_{gh}^{(N)}$ and $U_e^{(N)}=\Id$ at every length
    (Definition~\ref{def:mpug_group_representation}). The bond dimensions of
    $U_gU_h$ and $U_{gh}$ need not agree, so the product law holds for the
    periodic operators and not for the local tensors; comparing the stacked
    tensor $U_gU_h$ with $U_{gh}$ produces fusion tensors whose associativity
    defines a scalar three-cocycle
    (Section~\ref{sec:mpug_scalar_gauge}).
    % Source anchor: arXiv:2502.20257 (FrancoRubio2025MPUGauging),
    %   subsection "MPU group representations", as cited in
    %   def:mpug_group_representation; arXiv:2203.12563, lines 1026--1030.
\end{remark}

\begin{remark}[Group operators from a three-cocycle and the anomaly]%
    \label{rem:mposym_group_cocycle}
    \uses{def:mpug_three_cocycle, def:asymex_group_cocycle_tensor,
        thm:asymex_group_cocycle_mul, thm:asymex_group_cocycle_z2,
        thm:mpoalg_no_symmetric_normal}
    Every normalized three-cocycle $\omega$ of a finite group $G$
    (Definition~\ref{def:mpug_three_cocycle}) gives matrix product operators
    $U_g$ on the chain with site space $\C[G]$: a left shift by $g$ on every
    site, decorated by a diagonal two-site phase gate built from $\omega$
    (Definition~\ref{def:asymex_group_cocycle_tensor}). They satisfy
    $U_gU_h=U_{gh}$ at every length
    (Theorem~\ref{thm:asymex_group_cocycle_mul}). For the group of order two
    and its nontrivial cocycle, $U_s=\prod_iCZ_{i,i+1}Z_i\prod_iX_i$
    (Theorem~\ref{thm:asymex_group_cocycle_z2}), which is not a tensor
    product of one-site unitaries.

    The class of the associator in $H^3(G,\C^\times)$ is the
    \emph{anomaly}. An injective matrix product operator representation with
    $U_e=\Id$ whose fusion tensors reduce every stacked tensor $T_gT_h$
    exactly to $T_{gh}$ is on-site, and then its three-cocycle is trivial
    \cite[Section~4]{GarreRubio2023MPOSym}. Hence a nontrivial
    class forbids an on-site form. It also forbids an invariant normal matrix
    product state (Theorem~\ref{thm:mpoalg_no_symmetric_normal}). The sources
    argue further that a symmetric local Hamiltonian then cannot have a
    normal matrix product state as its unique ground state, and that such
    Hamiltonians are critical or symmetry breaking
    \cite{ChenLiuWen2011CZX,GarreRubio2023MPOSym,Cirac2021Matrix}.
    % Source anchor: arXiv:2203.12563, lines 677--681 (three-cocycle and
    %   H^3), 738 (no injective invariant state), 1211--1216 (injective
    %   representation with closed fusion is on-site), 2064--2070 (the
    %   cocycle construction and the Z_2 example); arXiv:2011.12127,
    %   Papers/2011.12127/TN-Review-main.tex, lines 1397--1419.
    % Lean: the on-site lemma of lines 1211--1216 is not formalized; planned
    %   name MPOTensor.GroupFamily.isOnSite_of_closed_fusion.
\end{remark}

%% ---------------------------------------------------------
\section{Fusion algebras of matrix product operators}%
\label{sec:mposym_fusion}

An invertible on-site symmetry acts on a matrix product state through a group
representation on the physical index. A matrix product operator symmetry is
given instead by a finite family of operators whose periodic operators
multiply according to fusion rules~\cite{GarreRubio2023MPOSym}. The source
defines the symmetry by invariance of the states with arbitrary boundary
conditions; the statements below use only its consequence for periodic
boundary conditions~\cite[Section~3]{GarreRubio2023MPOSym}; see
\cite{gap:glm23_mpo_symmetric_mps_scope}.

\begin{definition}[Matrix product operator fusion algebra]%
    \label{def:mpo_fusion_algebra}
    \lean{MPOTensor.IsMPOFusionAlgebra}
    \leanok
    \uses{def:mpo_tensor, def:mpo_family}
    A finite family $(T_a)_{a\in I}$ of matrix product operator tensors with
    structure constants $N_{ab}^c\in\N$ is a \emph{fusion algebra} if, for
    all $a,b\in I$ and $N\geq1$,
    \begin{align}
        O_a^{(N)}O_b^{(N)}
         & =\sum_{c\in I}N_{ab}^c\,O_c^{(N)}.
        \label{eq:mpo_fusion_algebra}
    \end{align}
    This is the fusion rule of \cite[Section~2]{GarreRubio2023MPOSym} and
    \cite{Bultinck2017Anyons}.
    % Source anchor: arXiv:2203.12563, lines 361--362.
\end{definition}

\begin{definition}[Unit and invertible labels]%
    \label{def:fusion_unit_invertible}
    \lean{MPOTensor.IsFusionUnit,
        MPOTensor.IsInvertibleLabel}
    \leanok
    \uses{def:mpo_fusion_algebra}
    A label $e$ is a \emph{unit} of the structure constants if
    $N_{eb}^c=N_{be}^c=\delta_{bc}$ for all $b,c$. A label $a$ is
    \emph{invertible} with respect to $e$ if there is a label $b$ with
    $N_{ab}^c=N_{ba}^c=\delta_{ce}$ for all $c$. The operator $O_e^{(N)}$ need
    not be the identity; outside the on-site case it is a projector
    \cite[Section~4]{GarreRubio2023MPOSym}.
    % Source anchor: arXiv:2203.12563, line 660.
\end{definition}

\begin{definition}[Fusion matrices and Perron--Frobenius dimension]%
    \label{def:mposym_pf_dimension}
    \lean{MPOTensor.fusionMatrix, MPOTensor.perronFrobeniusDim}
    \leanok
    \uses{def:mpo_fusion_algebra}
    Let $N_{ab}^c\in\N$ be structure constants on a finite label set $I$. The
    \emph{fusion matrix} of a label $a$ is the nonnegative integer matrix
    $N_a$ with entries $(N_a)_{bc}=N_{ab}^c$ for $b,c\in I$. The
    \emph{Perron--Frobenius dimension} of $a$ is the spectral radius
    $d_a=\rho(N_a)$.
    % Source anchor: arXiv:2204.05940, References/2204.05940/source/mpo.tex,
    %   line 6068 (N_a with (N_a)_b^c = N_{ab}^c, and d_a its spectral
    %   radius, the Perron--Frobenius dimension); arXiv:2203.12563,
    %   Papers/2203.12563/REsubmission.tex, line 1805 (r_a = d_a).
\end{definition}

\begin{theorem}[A nonnegative eigenvector for the Perron--Frobenius dimension]%
    \label{thm:mposym_pf_eigenvector}
    \lean{MPOTensor.exists_nonneg_eigenvector_perronFrobeniusDim,
        MPOTensor.perronFrobeniusDim_mem_spectrum,
        Matrix.exists_nonneg_mulVec_eq_spectralRadius_smul}
    \leanok
    \uses{def:mposym_pf_dimension}
    For every label $a$, the Perron--Frobenius dimension $d_a$ is an
    eigenvalue of $N_a$, and there is a column vector $v\neq0$ with
    nonnegative entries such that $N_av=d_av$.
\end{theorem}

\begin{proof}
    \leanok
    This is the Perron--Frobenius theorem for nonnegative matrices; we prove
    it for any square matrix $M$ with nonnegative real entries, and write
    $\rho=\rho(M)$. Choose a complex eigenvalue $\mu$ with $|\mu|=\rho$ and
    an eigenvector $u\neq0$, and let $|u|$ be the vector of moduli
    $|u_i|$. The triangle inequality gives $M|u|\geq\rho\,|u|$ entrywise.
    Hence the set $K$ of vectors $x$ with nonnegative entries,
    $\sum_ix_i=1$ and $Mx\geq\rho x$ entrywise is nonempty; it is compact
    and convex. For $x\in K$ the vector $(M+1)x$ is nonnegative with
    positive entry sum, and the map
    $x\mapsto(M+1)x/\sum_i((M+1)x)_i$ sends $K$ into itself, because
    $M(M+1)x-\rho(M+1)x=(M+1)(Mx-\rho x)$ and $M+1$ has nonnegative
    entries. The map is continuous on $K$. By
    Brouwer's fixed-point theorem it has a fixed point $x$, so $Mx=sx$ for
    some real $s$, and $x\neq0$ is nonnegative. From $x\in K$ we get
    $s\geq\rho$, and $s$ is an eigenvalue of $M$, so $s\leq\rho$. Thus
    $Mx=\rho x$ and $\rho$ is an eigenvalue. Applied to $M=N_a$ this gives
    the statement.
\end{proof}

\begin{lemma}[A positive left eigenvector gives the Perron--Frobenius dimension]%
    \label{lem:mposym_pf_dimension_positive}
    \lean{MPOTensor.perronFrobeniusDim_eq_of_pos_left_eigenvector,
        MPOTensor.perronFrobeniusDim_eq_of_dimension_relation}
    \leanok
    \uses{def:mposym_pf_dimension}
    Let $a$ be a label, $\lambda\in\R$, and let $(\delta_b)_{b\in I}$ be real
    numbers with $\delta_b>0$ for all $b$ and
    $\sum_b\delta_bN_{ab}^c=\lambda\,\delta_c$ for all $c$. Then $\lambda$ is
    the Perron--Frobenius dimension $d_a$. In particular, if positive numbers
    $\delta_b$ satisfy
    \begin{align}
        \sum_b N_{ab}^c\,\delta_b & =\delta_a\,\delta_c
        \label{eq:mposym_dimension_relation}
    \end{align}
    for all $a$ and $c$, then $\delta_a=d_a$ for every label $a$.
    % Source anchor: arXiv:2204.05940, mpo.tex line 6068 ("d_b defines a
    %   positive eigenvector for the matrix N_a ... the corresponding
    %   eigenvalue, d_a, is the spectral radius of N_a").
\end{lemma}

\begin{proof}
    \leanok
    Write $\delta$ for the row vector with entries $\delta_b$, so that
    $\delta N_a=\lambda\delta$. Since $\delta\neq0$, the matrix
    $\lambda-N_a$ has a nonzero left null vector, so its determinant
    vanishes and $\lambda$ is an eigenvalue of $N_a$. Let $\mu$ be any
    complex eigenvalue of $N_a$ with eigenvector $v\neq0$, and let $|v|$ be
    the vector of moduli $|v_b|$. The triangle inequality gives
    $|\mu|\,|v|\leq N_a|v|$ entrywise, and pairing with $\delta$ gives
    $|\mu|\,\delta|v|\leq\delta N_a|v|=\lambda\,\delta|v|$. Here
    $\delta|v|>0$ because every $\delta_b$ is positive and $v\neq0$, so
    $|\mu|\leq\lambda$. Thus $\lambda$ is an eigenvalue of largest modulus,
    and it is nonnegative, so $\lambda=\rho(N_a)=d_a$. The second statement
    is the first with $\lambda=\delta_a$.
\end{proof}

For matrix product operator representations of a weak Hopf algebra whose
unit corresponds to a single irreducible representation, the positive
dimensions attached to the labels satisfy
(\ref{eq:mposym_dimension_relation})
\cite[Section~5]{Molnar2022MPOAlgebrasI}, so by
Lemma~\ref{lem:mposym_pf_dimension_positive} they are the Perron--Frobenius
dimensions. For the Fibonacci and Ising rules below they are the quantum
dimensions of the corresponding fusion categories.
% Source anchor: arXiv:2204.05940, mpo.tex line 6068.

%% ---------------------------------------------------------
\section{Non-invertible symmetries}%
\label{sec:mposym_noninvertible}

A label without an inverse in the sense of
Definition~\ref{def:fusion_unit_invertible} gives an operator that has no
inverse within the fusion algebra. In the examples below such operators have
a nonzero kernel, so they are not invertible on the chain at all; they are
the lattice form of the non-invertible symmetries of Ising and anyon chains
\cite{Aasen2016Topological,Bultinck2017Anyons,Cirac2021Matrix}.

\begin{remark}[The unit is a projector]%
    \label{rem:mposym_unit_projector}
    \uses{def:fusion_unit_invertible, lem:asymex_fib_vacuum,
        thm:asymex_ising_path_operators}
    If $e$ is a unit of a fusion algebra, then $N_{ee}^c=\delta_{ec}$ gives
    $\bigl(O_e^{(N)}\bigr)^2=O_e^{(N)}$ and $O_e^{(N)}O_a^{(N)}=O_a^{(N)}
        O_e^{(N)}=O_a^{(N)}$ for every $a$. Thus $O_e^{(N)}$ is an idempotent
    and the algebra acts on its range; it is the identity only in special
    cases such as on-site symmetries
    \cite[Section~4]{GarreRubio2023MPOSym}. For the Fibonacci operators the
    unit is the orthogonal projection onto the admissible configurations
    (Lemma~\ref{lem:asymex_fib_vacuum}), and for the Ising operators it is
    the orthogonal projection onto the closed fusion paths
    (Theorem~\ref{thm:asymex_ising_path_operators}).
    % Source anchor: arXiv:2203.12563, line 660 (O_e is in general a
    %   projector rather than the identity).
\end{remark}

\begin{example}[Fibonacci fusion rule and the golden ratio]%
    \label{ex:mposym_fibonacci_dimension}
    \lean{FibonacciCompression.fusionMatrix_fibNim_tau,
        FibonacciCompression.charpoly_fusionMatrix_fibNim_tau,
        FibonacciCompression.spectrum_fusionMatrix_fibNim_tau,
        FibonacciCompression.fibNim_dimension_relation,
        FibonacciCompression.perronFrobeniusDim_fibNim_tau,
        FibonacciCompression.perronFrobeniusDim_fibNim_one,
        FibonacciCompression.not_isInvertibleLabel_fibNim_tau}
    \leanok
    \uses{def:mposym_pf_dimension, lem:mposym_pf_dimension_positive,
        def:asymex_fib_anomaly_nim, thm:asymex_fib_anomaly_fusion_algebra,
        def:asymex_fib_fsymbols}
    The Fibonacci labels are $1$ and $\tau$, with unit $1$ and
    $\tau\times\tau=1+\tau$. Its fusion matrix $N_\tau$ in the order
    $(1,\tau)$ is (\ref{eq:asymex_fib_anomaly_matrix}),
    with characteristic polynomial $\lambda^2-\lambda-1$ and eigenvalues
    $\varphi=(1+\sqrt5)/2$ and $-\varphi^{-1}$. Hence
    $d_\tau=\varphi$, the golden ratio, and $d_1=1$; the row vector
    $(1,\varphi)$ satisfies (\ref{eq:mposym_dimension_relation}). The label
    $\tau$ is not invertible: $\tau\times b$ contains $1$ only for
    $b=\tau$, and $\tau\times\tau$ also contains $\tau$. The operators of
    Section~\ref{sec:asymex_fibonacci} realize this rule at every length
    (Theorem~\ref{thm:asymex_fib_anomaly_fusion_algebra}), and the quantum
    dimensions entering their F-symbols are those of
    Definition~\ref{def:asymex_fib_fsymbols}
    \cite[Appendix~D.1]{Bultinck2017Anyons}.
    % Source anchor: arXiv:1511.08090, AnyonsPEPS.tex, lines 1240--1244
    %   (fusion rules, d_1 = 1, d_tau = phi).
\end{example}

\begin{example}[Ising fusion rule]%
    \label{ex:mposym_ising_dimension}
    \notready
    % Planned Lean: IsingTwist.perronFrobeniusDim_isingFusion_sigma.
    \uses{def:mposym_pf_dimension, lem:mposym_pf_dimension_positive,
        thm:asymex_ising_sigma_fusion, thm:asymex_ising_sigma_square}
    The Ising labels are $1$, $\psi$ and $\sigma$, with unit $1$ and the
    fusion rules of Section~\ref{sec:asymex_ising_tensors}, of which the only
    non-invertible one is $\sigma\times\sigma=1+\psi$. In the order
    $(1,\psi,\sigma)$ the fusion matrices are
    \begin{align}
        N_\psi   & =\begin{pmatrix}0&1&0\\1&0&0\\0&0&1\end{pmatrix}, &
        N_\sigma & =\begin{pmatrix}0&0&1\\0&0&1\\1&1&0\end{pmatrix}.
        \label{eq:mposym_ising_matrices}
    \end{align}
    The label $\psi$ is invertible and $N_\psi$ is a permutation matrix, so
    $d_\psi=1$. The eigenvalues of $N_\sigma$ are $0$ and $\pm\sqrt2$, so
    $d_\sigma=\sqrt2$, and $(1,1,\sqrt2)$ satisfies
    (\ref{eq:mposym_dimension_relation}). The operators of
    Section~\ref{sec:asymex_ising_tensors} realize these rules at every length
    (Theorems~\ref{thm:asymex_ising_sigma_fusion}
    and~\ref{thm:asymex_ising_sigma_square}). The relation
    $O_\sigma O_\psi=O_\sigma$ gives $O_\sigma(O_1-O_\psi)=0$, so $O_\sigma$
    annihilates the range of $O_1-O_\psi$.
    % Source anchor: arXiv:1511.08090, AnyonsPEPS.tex, lines 1306--1312
    %   (fusion rules and d_sigma = sqrt 2); arXiv:1601.07185,
    %   References/1601.07185/source/Ising-Defects.tex, lines 1051--1055.
\end{example}

\begin{remark}[Kramers--Wannier duality]%
    \label{rem:mposym_kramers_wannier}
    \uses{def:asymex_kw, cor:asymex_kw_square, lem:asymex_kw_sectors,
        rem:mpoalg_kramers_wannier}
    The Kramers--Wannier duality of the transverse-field Ising chain is a
    matrix product operator of bond dimension two
    (Definition~\ref{def:asymex_kw}). With the normalization
    $\mathcal D_N=2^{-N/2}K_N$ of Section~\ref{sec:asymex_kw}, its square is
    $\mathcal D_N^2=(\Id+\eta)T$, where $\eta$ is the global spin flip and $T$
    the one-site translation (Corollary~\ref{cor:asymex_kw_square} and
    (\ref{eq:asymex_kw_normalized_square})), and its kernel is the odd sector
    of $\eta$ (Lemma~\ref{lem:asymex_kw_sectors}). Aasen, Mong and Fendley
    keep the dual lattice distinct from the original one and obtain the Ising
    fusion rule $\mathcal D_\sigma^2=1+\mathcal D_\psi$, with the spin flip
    as $\mathcal D_\psi$ \cite{Aasen2016Topological}. On one periodic chain,
    identifying the two lattices produces the translation $T$, so the duality
    operator is not a member of a fusion algebra of operators on one chain.
    It is an intertwiner between the Ising chain and its dual
    \cite{Lootens2023Dualities}, as discussed in
    Remark~\ref{rem:mpoalg_kramers_wannier}.
    % Source anchor: arXiv:1601.07185, Ising-Defects.tex, lines 567 (duality
    %   commutes with the transfer matrix but is not invertible) and
    %   1040--1055 (D_sigma^4 = 2 D_sigma^2 and the fusion algebra);
    %   arXiv:2112.09091, References/2112.09091/source/_Examples.tex,
    %   lines 55--63 (the duality as a matrix product operator intertwiner);
    %   arXiv:2011.12127, TN-Review-main.tex, line 1397.
\end{remark}

%% ---------------------------------------------------------
\section{Symmetric matrix product states}%
\label{sec:mposym_symmetric}

A matrix product state is symmetric when each operator of the algebra maps
its periodic vectors into their span with length-independent coefficients.
For a single normal tensor the coefficients are eigenvalues; for several
normal tensors they form matrices. The results of this section show that
these coefficients are nonnegative integers compatible with the fusion rules,
which already excludes symmetric normal states for the Fibonacci rule. The
integrality rests on a property of normal tensors: a length-independent
relation between word traces has nonnegative integer coefficients.

\begin{definition}[Symmetric family of matrix product states]%
    \label{def:mpo_symmetric_family}
    \lean{MPOTensor.IsMPOSymmetricFamily,
        MPOTensor.IsMPOSymmetric}
    \leanok
    \uses{def:mpo_fusion_algebra, def:mps_tensor}
    A finite family $(A_x)_{x\in K}$ of matrix product state tensors is
    \emph{symmetric} under $(T_a)_{a\in I}$ with coefficients
    $M_{a,x}^y\in\C$ if, for all $a$, $x$ and $N\geq1$,
    \begin{align}
        O_a^{(N)}\ket{V^{(N)}(A_x)}
         & =\sum_{y\in K}M_{a,x}^y\ket{V^{(N)}(A_y)}.
        \label{eq:mpo_symmetric_family}
    \end{align}
    A single tensor $A$ is symmetric with eigenvalues $c_a$ if
    $O_a^{(N)}\ket{V^{(N)}(A)}=c_a\ket{V^{(N)}(A)}$ for all $a$ and $N\geq1$.
    The coefficients do not depend on $N$.
    % Source anchor: arXiv:2203.12563, lines 567--568 (restated at line 1801).
\end{definition}

\begin{definition}[Fusion character and nonnegative integer representation]%
    \label{def:fusion_character_nimrep}
    \lean{MPOTensor.IsFusionCharacter,
        MPOTensor.IsNIMRep}
    \leanok
    \uses{def:fusion_unit_invertible}
    A \emph{fusion character} with unit $e$ is a function $m$ on labels with
    $m_e=1$ and $m_am_b=\sum_cN_{ab}^cm_c$. A \emph{nonnegative integer
        representation} on a finite set $K$ is a family of matrices
    $M_{a,x}^y\in\N$ with
    \begin{align}
        \sum_cN_{ab}^cM_{c,x}^y
         & =\sum_{z\in K}M_{b,x}^zM_{a,z}^y.
        \label{eq:mpo_nimrep}
    \end{align}
    % Source anchor: arXiv:2203.12563, lines 564--565 (the relation for the
    %   multiplicities); the fusion character is not used by the source.
\end{definition}

\begin{theorem}[Word-trace relations against a normal tensor are integral]%
    \label{thm:asymex_fib_anomaly_trace_integral}
    \lean{MPSTensor.exists_nat_eq_of_forall_trace_evalWord_eq_mul}
    \leanok
    \uses{def:normal}
    Let $B$ be a matrix product state tensor and $A$ a word-normal tensor
    of positive bond dimension, over the same alphabet. If a scalar
    $c\in\C$ satisfies $\tr(B^w)=c\,\tr(A^w)$ for every nonempty word
    $w$, then $c$ is a nonnegative integer.
\end{theorem}

\begin{proof}
    \leanok
    Normality expresses a diagonal matrix unit as
    $E_{jj}=\sum_ux_uA^u$, where all words $u$ have the same positive
    length $\ell$. Define $Q=\sum_ux_uB^u$. Expanding powers and applying
    the trace hypothesis to the concatenated words gives, for every
    $k\geq1$,
    \begin{align}
        \tr(Q^k)
         & =c\,\tr\left(\left(\sum_ux_uA^u\right)^k\right)
        =c\,\tr(E_{jj}^k)=c.
    \end{align}
    A constant trace of all positive powers of a fixed finite-dimensional
    matrix is a nonnegative integer. Indeed, the power sums of its
    nonzero eigenvalues, counted with algebraic multiplicity, are
    constant only when every nonzero eigenvalue is one. The constant
    is their total multiplicity.
\end{proof}

\begin{theorem}[Rank-one actions on a normal tensor are integral]%
    \label{thm:asymex_fib_anomaly_integral_action}
    \lean{MPOTensor.exists_nat_eq_of_mpo_mulVec_mpv_eq_smul}
    \leanok
    \uses{thm:asymex_fib_anomaly_trace_integral, cor:asym_action_tensors}
    Let $T$ be a matrix product operator tensor and $A$ a word-normal
    matrix product state tensor of positive bond dimension, with the
    same physical alphabet. Suppose there is a single scalar $c\in\C$
    such that, for every $N\geq1$,
    \begin{align}
        O_N(T)\ket{V^{(N)}(A)}
         & =c\ket{V^{(N)}(A)}.
        \label{eq:asymex_fib_anomaly_scalar_action}
    \end{align}
    Then $c$ is a nonnegative integer.
\end{theorem}

\begin{proof}
    \leanok
    Corollary~\ref{cor:asym_action_tensors} identifies the periodic
    vector of $T\cdot A$ with $O_N(T)\ket{V^{(N)}(A)}$.
    Taking components in~(\ref{eq:asymex_fib_anomaly_scalar_action}) gives,
    for every nonempty word $w$,
    \begin{align}
        \tr((T\cdot A)^w) & =c\,\tr(A^w).
        \notag
    \end{align}
    Apply Theorem~\ref{thm:asymex_fib_anomaly_trace_integral}.
\end{proof}

\begin{theorem}[Integrality of multi-block word-trace relations]%
    \label{thm:mpoalg_multiblock_trace_integral}
    \lean{MPSTensor.exists_nat_eq_of_forall_trace_evalWord_eq_sum_mul}
    \leanok
    \uses{def:asym_word_module, def:normal}
    Let $(A_y)_{y\in K}$ be normal tensors of positive bond dimension with
    pairwise non-isomorphic word modules, let $B$ be a tensor on the same
    physical alphabet, and let $c_y\in\C$ satisfy
    $\tr(B^w)=\sum_yc_y\tr(A_y^w)$ for every nonempty word $w$. Then every
    $c_y$ is a nonnegative integer.
\end{theorem}

\begin{proof}
    \leanok
    \uses{def:asym_word_module}
    The word modules of the $A_y$ are simple and pairwise non-isomorphic, so
    by density the words of one positive length $\ell$ act on
    $\bigoplus_y\C^{D_y}$ as the full block-diagonal algebra. Fix $y_0$ and an
    index $j$; there are coefficients $x_u$, over words $u$ of length $\ell$,
    with
    \begin{align}
        \sum_ux_uA_{y_0}^u & =E_{jj},            &
        \sum_ux_uA_y^u     & =0\quad(y\neq y_0).
        \label{eq:mpoalg_trace_integral_unit}
    \end{align}
    Define $Q=\sum_ux_uB^u$. Expanding powers and applying the trace
    relation to the concatenated nonempty words gives, for every $k\geq1$,
    \begin{align}
        \tr(Q^k)
         & =\sum_yc_y\,\tr\left(\left(\sum_ux_uA_y^u\right)^k\right)
        =c_{y_0}\,\tr(E_{jj}^k)=c_{y_0}.
        \label{eq:mpoalg_trace_integral_powers}
    \end{align}
    A constant trace of all positive powers of a fixed finite-dimensional
    matrix is a nonnegative integer: the power sums of its nonzero
    eigenvalues, counted with algebraic multiplicity, are constant only when
    every nonzero eigenvalue is one, and the constant is their total
    multiplicity. This is the argument of
    Theorem~\ref{thm:asymex_fib_anomaly_trace_integral} with several blocks.
\end{proof}

\begin{theorem}[The action coefficients form a nonnegative integer representation]%
    \label{thm:mpo_symmetric_nimrep}
    \lean{MPOTensor.exists_nat_eq_of_isMPOSymmetricFamily,
        MPOTensor.sum_fusion_mul_eq_sum_mul_of_isMPOSymmetricFamily,
        MPOTensor.exists_isNIMRep_of_isMPOSymmetricFamily}
    \leanok
    \uses{def:mpo_symmetric_family, def:fusion_character_nimrep, def:normal}
    Let $(A_x)_{x\in K}$ be normal tensors of positive bond dimension whose
    periodic vectors are linearly independent at some length $N_0\geq1$,
    symmetric under a fusion algebra with coefficients $M_{a,x}^y$. Then
    every $M_{a,x}^y$ is a nonnegative integer and (\ref{eq:mpo_nimrep})
    holds. This is the integrality and associativity of the action
    multiplicities in \cite[Section~3]{GarreRubio2023MPOSym}.
    % Source anchor: arXiv:2203.12563, line 460 and lines 564--565.
\end{theorem}

\begin{proof}
    \leanok
    \uses{thm:mpoalg_multiblock_trace_integral}
    Fix $a$ and $x$. The word traces of $T_a\cdot A_x$ are
    $\sum_yM_{a,x}^y\tr(A_y^w)$ on nonempty words $w$. The word modules of
    the $A_y$ are simple and pairwise non-isomorphic, since isomorphic word
    modules have equal periodic vectors. An element of the augmentation ideal
    of the word algebra acts on the block $y_0$ as a matrix unit and
    annihilates the other blocks; its powers have constant trace
    $M_{a,x}^{y_0}$ on $T_a\cdot A_x$, so $M_{a,x}^{y_0}$ is a nonnegative
    integer (Theorem~\ref{thm:mpoalg_multiblock_trace_integral}). Applying
    (\ref{eq:mpo_fusion_algebra}) and (\ref{eq:mpo_symmetric_family}) twice
    to $O_a^{(N_0)}O_b^{(N_0)}\ket{V^{(N_0)}(A_x)}$ and comparing
    coefficients of the independent vectors gives (\ref{eq:mpo_nimrep}).
\end{proof}

\begin{theorem}[A single symmetric normal tensor gives a fusion character]%
    \label{thm:mpo_symmetric_character}
    \lean{MPOTensor.exists_isFusionCharacter_of_isMPOSymmetric}
    \leanok
    \uses{def:mpo_symmetric_family, def:fusion_character_nimrep}
    If a normal tensor of positive bond dimension is symmetric under a fusion
    algebra with eigenvalues $c_a$ and $c_e=1$, then $c_a\in\N$ for every $a$
    and $a\mapsto c_a$ is a fusion character.
\end{theorem}

\begin{proof}
    \leanok
    \uses{thm:asymex_fib_anomaly_integral_action}
    Each $c_a$ is a nonnegative integer by
    Theorem~\ref{thm:asymex_fib_anomaly_integral_action}. Applying
    (\ref{eq:mpo_fusion_algebra}) to the nonzero periodic vector at an
    injective length gives $c_ac_b=\sum_cN_{ab}^cc_c$.
\end{proof}

\begin{corollary}[Rank-one obstruction]%
    \label{cor:mpo_symmetric_no_go}
    \lean{MPOTensor.not_isMPOSymmetric_of_forall_not_isFusionCharacter}
    \leanok
    \uses{thm:mpo_symmetric_character}
    If the structure constants admit no fusion character with values in
    $\N$, no normal tensor of positive bond dimension is symmetric with
    $c_e=1$. This obstruction sees only the fusion ring; the cohomological
    obstruction of Theorem~\ref{thm:mpoalg_no_symmetric_normal}, which also
    excludes group-like algebras with a nontrivial three-cocycle, is a
    different statement.
\end{corollary}

\begin{proof}
    \leanok
    \uses{thm:mpo_symmetric_character}
    Theorem~\ref{thm:mpo_symmetric_character} produces a fusion character
    with values in $\N$ from any such tensor.
\end{proof}

\begin{theorem}[Invertible labels]%
    \label{thm:mpo_invertible_labels}
    \lean{MPOTensor.IsFusionCharacter.eq_one_of_isInvertibleLabel,
        MPOTensor.mpo_mulVec_eq_of_isInvertibleLabel,
        MPOTensor.IsNIMRep.exists_eq_one_of_isInvertibleLabel,
        MPOTensor.IsNIMRep.exists_equiv_of_isInvertibleLabel,
        MPOTensor.exists_equiv_mpo_mulVec_eq_of_isInvertibleLabel}
    \leanok
    \uses{thm:mpo_symmetric_nimrep, thm:mpo_symmetric_character,
        def:fusion_unit_invertible}
    A fusion character with values in $\N$ is $1$ on every invertible label,
    so an invertible label fixes the periodic vectors of a single symmetric
    normal tensor with $c_e=1$. For a family as in
    Theorem~\ref{thm:mpo_symmetric_nimrep} on which $O_e^{(N)}$ fixes every
    periodic vector, each invertible label $g$ permutes the blocks: there is a
    bijection $x\mapsto g\cdot x$ of the block labels with
    $O_g^{(N)}\ket{V^{(N)}(A_x)}=\ket{V^{(N)}(A_{g\cdot x})}$ for all
    $N\geq1$. This is the permutation action of
    \cite[Section~4]{GarreRubio2023MPOSym}.
    % Source anchor: arXiv:2203.12563, line 683.
\end{theorem}

\begin{proof}
    \leanok
    \uses{thm:mpo_symmetric_nimrep, thm:mpo_symmetric_character}
    From $m_am_b=m_e=1$ in $\N$ one gets $m_a=1$. For the family, the
    multiplicity matrices satisfy $M_aM_b=M_bM_a=M_e=\Id$ with nonnegative
    integer entries, so each row of $M_a$ is a standard basis vector, and the
    two relations make the resulting map of blocks a bijection.
\end{proof}

\begin{example}[Fibonacci]%
    \label{ex:mpo_symmetry_fibonacci}
    \lean{FibonacciCompression.isMPOFusionAlgebra_fibBlock,
        FibonacciCompression.isFusionUnit_fibFusion,
        FibonacciCompression.isMPOSymmetricFamily_fibNimTargets,
        FibonacciCompression.isNIMRep_fibFusion,
        FibonacciCompression.not_isFusionCharacter_fibFusion}
    \leanok
    \uses{cor:mpo_symmetric_no_go, thm:asymex_fib_anomaly_fusion_algebra,
        thm:asymex_fib_anomaly_nim_rep, thm:asymex_fib_anomaly_no_go}
    The Fibonacci operators of Section~\ref{sec:asymex_fibonacci} form a
    fusion algebra with unit $1$ and $\tau\times\tau=1+\tau$; the two normal
    states with $\tau\cdot x_1=x_\tau$ and $\tau\cdot x_\tau=x_1+x_\tau$ form
    a symmetric family whose action is the regular representation
    \cite[Section~6]{GarreRubio2023MPOSym}. Since $m^2=1+m$ has no solution
    in $\N$, no single normal tensor $A$ satisfies
    $O_1^{(N)}\ket{V^{(N)}(A)}=\ket{V^{(N)}(A)}$ and
    $O_\tau^{(N)}\ket{V^{(N)}(A)}=c\ket{V^{(N)}(A)}$ for all $N\geq1$
    (Theorem~\ref{thm:asymex_fib_anomaly_no_go}).
    % Source anchor: arXiv:2203.12563, line 1993; arXiv:1511.08090, App. D.1.
\end{example}

\begin{example}[Anomalous cyclic group of order three]%
    \label{ex:mpo_symmetry_z3}
    \lean{Z3Anomalous.isMPOFusionAlgebra_z3Block,
        Z3Anomalous.isInvertibleLabel_z3Fusion}
    \leanok
    \uses{def:fusion_unit_invertible}
    The operators $1,U,U^\dagger$ of
    Section~\ref{sec:asymex_z3anomalous} form a group-like fusion algebra in
    which every label is invertible. The fusion ring does not see the
    nontrivial three-cocycle of this representation; the obstruction to
    invariant normal states comes from
    Theorem~\ref{thm:mpoalg_no_symmetric_normal}.
\end{example}

\begin{remark}[Symmetric superpositions of several normal states]%
    \label{rem:mposym_pf_symmetric_state}
    \uses{thm:mpo_symmetric_nimrep, ex:mpo_symmetry_fibonacci,
        def:mposym_pf_dimension}
    The obstruction of Corollary~\ref{cor:mpo_symmetric_no_go} concerns a
    single normal tensor. A superposition of the periodic vectors of several
    normal tensors can still be an eigenvector of every operator. If positive
    numbers $v_x$ satisfy $\sum_xv_xM_{a,x}^y=r_av_y$ for all $a$ and $y$,
    with the coefficients $M_{a,x}^y$ of
    Theorem~\ref{thm:mpo_symmetric_nimrep}, then
    $\ket{\psi^{(N)}}=\sum_xv_x\ket{V^{(N)}(A_x)}$ satisfies
    $O_a^{(N)}\ket{\psi^{(N)}}=r_a\ket{\psi^{(N)}}$; Garre-Rubio, Lootens
    and Molnár show that such a $v$ exists when the algebra is that of a
    fusion category \cite[Section~5]{GarreRubio2023MPOSym}. For the Fibonacci
    family of Example~\ref{ex:mpo_symmetry_fibonacci}, $M_\tau=N_\tau$ and
    $v=(1,\varphi)$, so
    $\ket{V^{(N)}(A_1)}+\varphi\ket{V^{(N)}(A_\tau)}$ is an eigenvector of
    $O_\tau^{(N)}$ with eigenvalue $d_\tau=\varphi$, which is not an integer.
    % Source anchor: arXiv:2203.12563, lines 1795--1805 (common positive
    %   eigenvector of the M_a, eigenvalue r_a, and r_a = d_a when
    %   M_a = N_a) and line 1993 (the Fibonacci action).
\end{remark}

%% ---------------------------------------------------------
\section{The zipper condition}%
\label{sec:mposym_zipper}

\begin{remark}[Zipper and pulling-through conditions]%
    \label{rem:mposym_zipper}
    \uses{def:mpdo_complete_zipper_fusion_family, rem:asym_zipper,
        def:mpoalg_fusion_tensors, def:mpoalg_action_tensors,
        thm:mpoalg_zipper_fusion_uniqueness}
    The fusion rule (\ref{eq:mpo_fusion_algebra}) is an identity of periodic
    operators. Its local form, the \emph{zipper condition}, asks for maps
    $X_{ab,\mu}^c\colon\C^{\chi_c}\to\C^{\chi_a}\otimes\C^{\chi_b}$ and left
    inverses $X_{ab,\mu}^{c+}$, for $\mu=1,\ldots,N_{ab}^c$, with
    $X_{ab,\nu}^{d+}X_{ab,\mu}^c=\delta_{cd}\delta_{\mu\nu}\Id_{\chi_c}$ and
    \begin{align}
        (T_aT_b)^{ik}X_{ab,\mu}^c    & =X_{ab,\mu}^c\,T_c^{ik},
        \notag                                                    \\
        X_{ab,\mu}^{c+}(T_aT_b)^{ik} & =T_c^{ik}\,X_{ab,\mu}^{c+}
        \label{eq:mposym_zipper}
    \end{align}
    for all physical indices $i,k$ \cite{Bultinck2017Anyons}. A fusion
    tensor can then be slid through the product of two operators one site at
    a time. If in addition the product is reconstructed on the span of the zipper
    maps, $(T_aT_b)^{ik}=\sum_{c,\mu}X_{ab,\mu}^cT_c^{ik}X_{ab,\mu}^{c+}$,
    as required in Definition~\ref{def:mpdo_complete_zipper_fusion_family},
    then closing the chain gives (\ref{eq:mpo_fusion_algebra}) at every
    length. Its pictorial form, in which an operator is pulled through
    a region of a two-dimensional tensor network, is the
    \emph{pulling-through} condition of topologically ordered projected
    entangled pair states \cite{Bultinck2017Anyons,Cirac2021Matrix}. The
    zipper identities appear as (\ref{eq:rfp_zipper_synthesis}) and
    (\ref{eq:rfp_zipper_analysis}) in
    Definition~\ref{def:mpdo_complete_zipper_fusion_family}.

    The converse fails: equality of periodic operators does not imply a
    sitewise intertwiner (Remark~\ref{rem:asym_zipper}), and the group
    operators of the nontrivial three-cocycle of the group of order two
    satisfy only one of the two identities
    \cite[Section~6]{GarreRubio2023MPOSym}. Chapter~\ref{ch:mpo_algebras}
    therefore obtains fusion and action tensors from equality of periodic
    traces through reductions
    (Definitions~\ref{def:mpoalg_fusion_tensors}
    and~\ref{def:mpoalg_action_tensors}), and shows that fusion tensors
    with vanishing remainder, the situation of the zipper condition, are
    unique up to invertible multiplicity matrices
    (Theorem~\ref{thm:mpoalg_zipper_fusion_uniqueness}).
    % Source anchor: arXiv:1511.08090, AnyonsPEPS.tex, lines 159--170
    %   (fusion tensors and left inverses), 194--200 (the zipper condition),
    %   407--419 (pulling through); arXiv:2011.12127, TN-Review-main.tex,
    %   lines 1311--1330 (zipper and pulling-through equations);
    %   arXiv:2203.12563, lines 2224--2243 (left zipper only for the Z_2
    %   group operator of the nontrivial cocycle).
\end{remark}
